% Unpublished student project draft -- NOT submitted anywhere, NOT peer-reviewed.
% Content mirrors research/paper/draft.md (the authoritative source); this is a
% structural LaTeX skeleton, not a second independently-written draft, so the
% two never drift into contradicting each other.
\documentclass[conference]{IEEEtran}
\usepackage{cite}
\usepackage{amsmath}
\usepackage{graphicx}
\usepackage{booktabs}

\begin{document}

\title{Evidence-Adaptive Retrieval for OCR-Resilient Document Question Answering}
\author{\IEEEauthorblockN{Manthira Moorthy P, Vijay S K}
\IEEEauthorblockA{IV ECE-A \\ Mentor: Dr. V. Karthikeyan, AP(Sl.G.)}}

\maketitle

\begin{abstract}
Retrieval-augmented QA systems, including the QuIM-RAG baseline this project extends
(Saha et al., IEEE Access 2024), retrieve top-$k$ passages and generate an answer
unconditionally regardless of retrieval confidence. We extend an existing OCR-aware
hybrid-retrieval prototype with real BM25 lexical scoring, real per-chunk OCR
confidence, a 2-stage residual vector quantizer, and Evidence-Adaptive Retrieval \&
Verification (EAR): a logistic combination of four zero-extra-inference signals --
dense/BM25 agreement, OCR confidence, cross-passage numeric contradiction, and query
complexity -- calibrated on held-out data to decide whether to answer, retrieve more,
inspect structure, or abstain. On a stratified 420-question subset of OHR-Bench,
restricting answers to EAR's confident decisions raises top-1 retrieval accuracy from
71.2\% to 78.3\% (62.8\% coverage; 95\% bootstrap CI [+2.5, +12.0] points). We also
report, without softening, two findings that complicate a simple ``our method wins''
narrative: real BM25 alone outperforms the existing hybrid fusion on this benchmark,
and our RVQ implementation's compression gains come at a Recall@5 cost production
ANN systems avoid via techniques not implemented here.
\end{abstract}

\section{Introduction and Motivation}
% See research/PROJECT_AUDIT.md and research/paper/draft.md Section 1 for full text.
Real-world document QA over scanned, OCR'd, and noisy documents is where hybrid
retrieval is most likely to retrieve the wrong evidence, yet most RAG pipelines
answer regardless of retrieval confidence.

\section{Related Work}
% Full annotated table: research/PAPER_REVIEW.md
The closest published baseline is Saha et al.~\cite{saha2024quimrag}, evaluated on a
private corpus with no OCR noise, no lexical retrieval path, and no abstention
mechanism. Recent work on RAG abstention and retrieval-agreement-as-signal is
surveyed in \texttt{research/PATENTABILITY\_ASSESSMENT.md}.

\section{Research Gap}
% research/RESEARCH_GAP_ANALYSIS.md

\section{Problem Formulation}
Given query $q$ and evidence $E_k(q)$ retrieved for it, we seek a decision function
$f(q, E_k(q)) \rightarrow \{\text{answer}, \text{expand}, \text{inspect}, \text{abstain}\}$
that improves expected correctness of answered queries over the unfiltered baseline.

\section{Proposed Method: Evidence-Adaptive Retrieval (EAR)}
% Full equations and pseudocode: research/PROPOSED_CONTRIBUTION.md
\begin{equation}
p = \sigma(w_0 + w_1 a + w_2 o + w_3 x + w_4 c)
\end{equation}
where $a$ = dense/BM25 rank-agreement, $o$ = OCR confidence, $x$ = contradiction, $c$
= query complexity, and $w$ is fit by logistic regression on a calibration split.

\section{System Architecture}
% Diagram source: research/paper/draft.md Section 6 (ASCII pipeline diagram)
\begin{figure}[h]
\centering
% \includegraphics[width=\columnwidth]{architecture.png}
\caption{Ingestion through EAR decision to answer/expand/inspect/abstain. See
research/paper/draft.md Section 6 for the pipeline diagram (not yet rendered as a
standalone figure file).}
\end{figure}

\section{Experimental Setup}
\texttt{research/run\_benchmark.py} against \texttt{research/data/ohrbench\_adapter.py}.

\section{Baselines}
Dense-only, real BM25-only, Hybrid RRF, Hybrid+rerank, Hybrid+rewriting,
Hybrid+rewriting+rerank.

\section{Dataset}
OHR-Bench~\cite{zhang2024ohrbench}, a stratified 420-QA-pair subset, 6 domains, seed
42, 40/60 calibration/evaluation split.

\section{Evaluation Metrics}
% research/METRIC_INVENTORY.md

\section{Results and Statistical Analysis}
% Full tables: research/BASELINE_COMPARISON.md Sections 1-2
\begin{table}[h]
\centering
\caption{Baseline retrieval comparison (k=5, n=252)}
\begin{tabular}{lccc}
\toprule
Config & nDCG@5 & MRR@5 & $p$ vs.\ Dense \\
\midrule
Dense only & 0.477 & 0.726 & -- \\
BM25 (real) & 0.528 & 0.837 & 0.000 \\
Hybrid (RRF) & 0.515 & 0.780 & 0.000 \\
Hybrid + rerank & 0.574 & 0.870 & 0.000 \\
Hybrid + rewriting & 0.453 & 0.694 & 0.059 \\
Hybrid + rewriting + rerank & 0.573 & 0.875 & 0.000 \\
\bottomrule
\end{tabular}
\end{table}

EAR: 71.2\% $\rightarrow$ 78.3\% top-1 accuracy at 62.8\% coverage, 95\% CI
[+2.5, +12.0] points (5{,}000-resample bootstrap).

\section{Ablation Studies}
% research/BASELINE_COMPARISON.md Section 2 -- leave-one-signal-out, refit per ablation.

\section{Error Analysis}
Not performed at the per-query qualitative level in this build; see
\texttt{research/RESEARCH\_GAP\_ANALYSIS.md}.

\section{Computational Cost Analysis}
% research/BASELINE_COMPARISON.md latency columns

\section{Limitations}
% research/RESEARCH_GAP_ANALYSIS.md, inline notes in research/BASELINE_COMPARISON.md

\section{Threats to Validity}
Single benchmark, single seed, single embedding model; results may not transfer
without re-calibration.

\section{Reproducibility Statement}
\texttt{.venv/bin/python research/data/ohrbench\_adapter.py \&\& .venv/bin/python
research/run\_benchmark.py} (seed=42; \textasciitilde6 minutes on the development
machine). Raw logs: \texttt{research/experiments/20261014T185643Z.json}.

\section{Conclusion and Future Work}
A calibrated, zero-extra-inference sufficiency score measurably improves selective
answer accuracy, but the ablation shows this is driven mostly by one signal
(agreement), not uniformly by all four as hypothesized.

\bibliographystyle{IEEEtran}
\begin{thebibliography}{9}
\bibitem{saha2024quimrag} B. Saha, U. Saha, and M. Z. Malik, ``QuIM-RAG: Advancing
Retrieval-Augmented Generation With Inverted Question Matching for Enhanced QA
Performance,'' IEEE Access, vol. 12, pp. 185401--185410, 2024.
\bibitem{zhang2024ohrbench} J. Zhang et al., ``OCR Hinders RAG: Evaluating the
Cascading Impact of OCR on Retrieval-Augmented Generation,'' arXiv:2412.02592, 2024.
\bibitem{jegou2011pq} H. J\'egou, M. Douze, and C. Schmid, ``Product Quantization for
Nearest Neighbor Search,'' IEEE TPAMI, vol. 33, no. 1, pp. 117--128, 2011.
\bibitem{robertson2009bm25} S. Robertson and H. Zaragoza, ``The Probabilistic
Relevance Framework: BM25 and Beyond,'' Foundations and Trends in IR, vol. 3, no. 4,
pp. 333--389, 2009.
\end{thebibliography}

\end{document}
